% Section: Conclusions — P8 commit 9.

\section{Conclusions}
\label{sec:conclusions}

We pre-registered and executed a quantum PINN benchmark testing the
Schuld--Fourier regime-dependent advantage
hypothesis~\cite{schuld2022right,schuld2021effect}: four quantum-PINN
families on a hardness ladder of four ODEs and four PDEs at three seeds
each, against five matched classical baselines and the complete
2$\times$2 fairness matrix on the forecaster task.

\textbf{The pre-registered $H_1$ regime-dependent quantum
advantage is FALSIFIED.} Across all eight non-fragile sensitivity
points, the paired-bootstrap 95\% confidence intervals on
$\dsmooth - \dbroad$ either include zero or exclude it negatively. The
primary solver verdict at $n = 24$ has $\ddiff = -0.084$, CI
$[-0.278, +0.061]$; the primary forecaster verdict at $n = 9$ has
$\ddiff = -0.501$, CI $[-0.804, -0.244]$, excluding zero in the
direction opposite the prediction.

\textbf{The LTC decomposition attributes the inversion to the quantum
circuit.} The quantum-isolated $\ddiff = -0.616$ (CI $[-1.167, -0.178]$)
is significantly negative, while the liquid-isolated
$\ddiff = +0.115$ (CI $[-0.097, +0.377]$) is the only verdict whose
point estimate matches the Schuld--Fourier direction: the classical
liquid-time-constant machinery carries the predicted smooth-bin
inductive bias, and the quantum circuit destroys it.

\textbf{One positive sub-finding survives.} At zero classical
parameters in the quantum cell, \texttt{te\_qpinn\_qnn} on \fhn{}
attains 24$\times$ tighter seed standard deviation and 2.5$\times$ lower
mean error than the matched cPINN. The H3 cross-tabulation flags
expressivity (KL-to-Haar distance) as a candidate within-roster
predictor of per-cell $\Delta$ ($\rho = +0.518$, LOO-robust at
$n = 9$), underpowered but directionally consistent with the
decomposition's mechanism attribution.

\textbf{The methodological contribution is the framework itself.} The
open-source pipeline at \url{https://github.com/shawngibford/qlnn}
implements the Bowles--Schuld~\cite{bowles2024better} program end to
end, and its P7.11 ablation completes the forecaster fairness
2$\times$2 -- to our knowledge the first quantum-PINN benchmark to do so.

\textbf{Future work targets the mechanism.} The follow-up paper applies
three interventions the mechanism analysis identifies as relevant:
Quantum Natural Gradient optimization~\cite{stokes2020quantum} for the
quantum-circuit training inefficiency, causal training
weighting~\cite{wang2022causality} for the stiff-dynamics failure mode,
and data-reuploading depth $L = 5$ for the Fourier $K_{\max}$ ceiling
(Section~\ref{sec:mechanism-scalars}); the deferred Kuramoto and KdV
systems are queued alongside. Whether any recovers the $H_1$ direction
is an empirical question the framework introduced here is built to keep
falsifiable.
